\documentclass[12pt]{article}

% ---------------------------------------------------------------------
%  PACKAGES
% ---------------------------------------------------------------------
\usepackage{fontawesome}   % \faEnvelopeO icon in author footnotes
\usepackage{titling}       % \emptythanks (clears \thanks between title pages)
\usepackage{bibunits}      % separate bibliographies for main text and supplement
\usepackage{amsmath}
\usepackage{amsfonts}
\usepackage{natbib}        % \citet, \citep, \citeauthor, \citeyearpar
\usepackage{url}
\usepackage[breaklinks=true,bookmarksopen=true,colorlinks=true,citecolor=blue]{hyperref}

\usepackage{graphicx}      % \includegraphics
\usepackage{booktabs}      % \toprule, \midrule, \bottomrule
\usepackage{caption}       % \captionsetup
\usepackage{threeparttable}% tables with notes
\usepackage{float}         % [H] placement
\usepackage{setspace}      % \begin{spacing}{...}
\usepackage{xcolor}

% ---------------------------------------------------------------------
%  BLINDING SWITCH
%  1 = show title page with authors;  0 = blinded version (title only)
% ---------------------------------------------------------------------
\newcommand{\blind}{1}

% ---------------------------------------------------------------------
%  MARGINS (1 inch all around -- do not change)
% ---------------------------------------------------------------------
\addtolength{\oddsidemargin}{-.7in}%
\addtolength{\evensidemargin}{-.7in}%
\addtolength{\textwidth}{1.4in}%
\addtolength{\textheight}{1.3in}%
\addtolength{\topmargin}{-.8in}%

% ---------------------------------------------------------------------
%  REUSABLE TITLE AND AUTHOR BLOCK
%  Edit these once; they are used on the main and supplement title pages.
% ---------------------------------------------------------------------
\newcommand{\papertitle}{Paper Title}

\newcommand{\paperauthors}{%
    Author 1\thanks{Corresponding Author: Address 1. \faEnvelopeO: \href{mailto:author1@email.edu}{author1@email.edu}.} \and
    Author 2\thanks{Address 2. \faEnvelopeO: \href{mailto:author2@email.edu}{author2@email.edu}.} \and
    Author 3\thanks{Address 3. \faEnvelopeO: \href{mailto:author3@email.edu}{author3@email.edu}.} \and
    Author 4\thanks{Address 4. \faEnvelopeO: \href{mailto:author4@email.edu}{author4@email.edu}.}%
}

% =====================================================================
\begin{document}
\begin{bibunit}[jpe]

\def\spacingset#1{\renewcommand{\baselinestretch}%
{#1}\small\normalsize} \spacingset{1}

% ---------------------------------------------------------------------
%  TITLE PAGE
% ---------------------------------------------------------------------
\if1\blind
{
\title{\papertitle\thanks{Acknowledgments go here. Thank seminar participants, funding sources, and research assistants. State where replication data and code are available. All errors are our own.}}
\author{\paperauthors}
\maketitle
} \fi

\if0\blind
{
  \bigskip
  \bigskip
  \bigskip
  \begin{center}
    {\LARGE\bf \papertitle}
  \end{center}
  \medskip
} \fi

\bigskip
\begin{abstract}
\noindent Replace this text with the abstract. State the question, the data, the method, and the main finding in roughly 150--250 words. Cite as needed, e.g.\ \citet{entry1} or \citep{entry2}.
\end{abstract}

\noindent%
{\it Keywords:} Keyword 1; Keyword 2; Keyword 3.
\\
\noindent%
{\it JEL code:} C01, C13, J24.

\vfill

\newpage

% ---------------------------------------------------------------------
%  MAIN TEXT
% ---------------------------------------------------------------------
\spacingset{1.56} % DON'T change the spacing!

\section{Introduction}\label{sec:intro}
Motivate the question and summarize the contribution. Reference related work with natbib, for example \citet{entry1} in text, or in parentheses \citep{entry2, entry3}. Point the reader forward with cross-references such as Section~\ref{sec:data} and Section~\ref{sec:results}.

\section{Background}\label{sec:background}
Describe the institutional setting or the literature this paper builds on. Display equations are numbered:
\begin{equation}
    y_i = \alpha + \beta x_i + \varepsilon_i, \qquad i = 1, \ldots, n.
    \label{eq:linear}
\end{equation}
Equation~\eqref{eq:linear} can then be referenced in the text.

\href{}{}\section{Data}\label{sec:data}
Describe the data sources, the sample, and the key variables. Table~\ref{tab:summary} is a template for a summary-statistics table with notes.

\begin{table}[H]
\centering
\begin{threeparttable}
\caption{Summary Statistics}
\label{tab:summary}
\begin{tabular}{lcccc}
\toprule
Variable & Mean & Std.\ Dev. & Min & Max \\
\midrule
Variable 1 & 0.00 & 0.00 & 0.00 & 0.00 \\
Variable 2 & 0.00 & 0.00 & 0.00 & 0.00 \\
Variable 3 & 0.00 & 0.00 & 0.00 & 0.00 \\
\midrule
Observations & \multicolumn{4}{c}{0} \\
\bottomrule
\end{tabular}
\begin{tablenotes}[flushleft]
\footnotesize
\item \textit{Notes:} Describe the sample and the source of each variable here.
\end{tablenotes}
\end{threeparttable}
\end{table}

\section{Empirical Methodology}\label{sec:methodology}
Lay out the estimation strategy and identification argument. Multi-line derivations can use the \texttt{align} environment:
\begin{align}
    \mathbb{E}[y_i \mid x_i] &= g(x_i^{\prime}\beta), \label{eq:cef} \\
    \hat{\beta} &= \arg\max_{\beta} \sum_{i=1}^{n} \ell_i(\beta). \label{eq:mle}
\end{align}

\section{Results}\label{sec:results}
Report and interpret the estimates. Figure~\ref{fig:example} shows how to include a graphic; replace the placeholder box with \verb|\includegraphics[width=0.8\textwidth]{figures/filename.pdf}|.

\begin{figure}[H]
\centering
\fbox{\parbox[c][5cm][c]{0.8\textwidth}{\centering Figure placeholder}}
\caption{Example Figure}
\label{fig:example}
\begin{minipage}{0.9\textwidth}
\footnotesize \textit{Notes:} Describe what is plotted and the data source.
\end{minipage}
\end{figure}

\section{Discussion and Conclusion}\label{sec:conclusion}
Summarize the findings, discuss limitations, and suggest directions for future work. Additional material is collected in \ref{Appendix_A} and \ref{Appendix_B} of the Supplemental Materials.

\putbib[refs]
\end{bibunit}

% =====================================================================
%  SUPPLEMENTAL MATERIALS
%  Delete everything from here to \end{document} if not needed.
% =====================================================================
\emptythanks
\clearpage
\setcounter{page}{1}

\if1\blind
{
  \title{\papertitle\\ -- Supplemental Materials --}
  \author{\paperauthors}
  \maketitle
} \fi

\if0\blind
{
  \bigskip
  \bigskip
  \bigskip
  \title{\bf \papertitle\\ -- Supplemental Materials --}
  \maketitle
  \medskip
} \fi

% ---------------------------------------------------------------------
%  APPENDIX SETTINGS
% ---------------------------------------------------------------------
\setcounter{section}{0}
\appendix

% Section titles read "Appendix A", "Appendix B", ...
\renewcommand\thesection{Appendix \Alph{section}}

% Tables and figures numbered A-1, A-2, B-1, ... and reset in each appendix
\renewcommand{\thetable}{\Alph{section}-\arabic{table}}
\renewcommand{\thefigure}{\Alph{section}-\arabic{figure}}
\makeatletter
\@addtoreset{table}{section}
\@addtoreset{figure}{section}
\makeatother

% ---------------------------------------------------------------------
%  APPENDIX CONTENT
% ---------------------------------------------------------------------
\begin{bibunit}[jpe]
\begin{spacing}{1.45}

\noindent This is the supplementary material for Author 1, Author 2, Author 3 and Author 4 (``\papertitle,'' \the\year). Briefly describe what each appendix contains. \ref{Appendix_A} collects additional tables; \ref{Appendix_B} collects proofs and derivations.

\section{Additional Tables\label{Appendix_A}}
Appendix tables are numbered automatically as Table~\ref{tab:appA1}.

\begin{table}[H]
\centering
\begin{threeparttable}
\caption{Robustness Checks}
\label{tab:appA1}
\begin{tabular}{lcc}
\toprule
 & (1) & (2) \\
\midrule
Coefficient 1 & 0.000 & 0.000 \\
              & (0.000) & (0.000) \\
Coefficient 2 & 0.000 & 0.000 \\
              & (0.000) & (0.000) \\
\midrule
Observations & 0 & 0 \\
\bottomrule
\end{tabular}
\begin{tablenotes}[flushleft]
\footnotesize
\item \textit{Notes:} Standard errors in parentheses. See \citet{entry4} for the estimator.
\end{tablenotes}
\end{threeparttable}
\end{table}

\section{Proofs and Derivations\label{Appendix_B}}
Collect proofs here. Equations in the appendix continue the numbering of the main text unless you redefine \verb|\theequation|.

\putbib[refs]
\end{spacing}
\end{bibunit}

\end{document}
